%% trie -- the reverse trie of the running example at level 1.
\begin{figure}[t]
\centering
\begin{tikzpicture}[font=\scriptsize, x=1.5cm, y=1cm,
    nd/.style={circle, draw, fill=white, inner sep=1.3pt},
    en/.style={circle, draw, fill=black, inner sep=1.3pt},
    pt/.style={circle, fill=black!40, inner sep=0.7pt}]
%% depth d is drawn at x = -d; root at the right
\node[nd] (r) at (0,0) {};
%% branch ending with 0: parses 0, 010, 0010
\node[en] (a1) at (-1,0.6) {};
\node[pt] (a2) at (-2,0.6) {};
\node[en] (a3) at (-3,0.6) {};
\node[en] (a4) at (-4,0.6) {};
\draw (r) -- node[above right, inner sep=1.5pt] {$0$} (a1);
\draw (a1) -- node[below] {$1$} (a2);
\draw (a2) -- node[below] {$0$} (a3);
\draw (a3) -- node[below] {$0$} (a4);
\node[above=2pt] at (a1) {$0$};
\node[above=2pt] at (a3) {$010$};
\node[left=2pt] at (a4) {$Z=0010$};
%% branch ending with 1: parses 11, 011
\node[pt] (b1) at (-1,-0.6) {};
\node[en] (b2) at (-2,-0.6) {};
\node[en] (b3) at (-3,-0.6) {};
\draw (r) -- node[below right, inner sep=1.5pt] {$1$} (b1);
\draw (b1) -- node[above] {$1$} (b2);
\draw (b2) -- node[above] {$0$} (b3);
\node[below=2pt] at (b2) {$11$};
\node[left=2pt] at (b3) {$011$};
%% edges of the path of Z
\foreach \a/\b/\l in {0/1/{(Z,0,1)},1/3/{(Z,1,3)},3/4/{(Z,3,4)}} {
  \draw[decorate, decoration={brace, amplitude=3pt}] (-\b+0.06,1.35) -- (-\a-0.06,1.35)
    node[midway, above=3pt] {$\l$};
}
%% depth axis
\draw[->, >={Stealth[length=1.8mm]}] (0,-1.65) -- (-4.5,-1.65) node[left] {depth};
\foreach \d in {0,...,4} { \draw (-\d,-1.58) -- (-\d,-1.72) node[below] {$\d$}; }
\end{tikzpicture}
\caption{The reverse trie of \cref{ex:trie}: the parses of the five upper
blocks, read from the right; the edges of the drawing are labeled with
lower blocks, here letters. A black node is the end of the parse written
next to it. The path of
$Z=0010$ has the nodes $0,1,3,4$; nothing branches off at depth $2$, so
$(Z,1,3)$ is one edge, which is traversed by the upper blocks $0010$ and
$010$.}
\label{fig:trie}
\end{figure}
